\section{Discussion and limitations}
\label{sec:discussion}

\subsection{Correction-aware authority as a data relation}

The result links three roles that a final value cannot recover: proposal origin, human value authority, and committed field source. Admission binds these roles to one candidate and fresh predecessor while constructing a complete successor. The paired histories explain the corresponding observation distinctions; the formal, projection, and transaction checks test their enforcement.

The journal comparison locates the contribution in this relation. Exact binding gives the journal and reference the same 165 decisions and query-answer sets. Instance accountability requires rejecting the 15 context-equivalent substitutions; value/context accountability permits them. Applications should choose the granularity that their audit questions require. Renewed approval without a value change belongs to a separate review-event path, because concrete post-initial admission creates transitions only for changed values.

\subsection{Trust and validity boundaries}

The server-held gate demonstrates request-to-review consistency for the tested substitutions. Display-only controls retain a separate presentation boundary. Host authentication, review-channel integrity, principal policy, hashing/canonicalization, and database execution remain trusted. Recorded lineage identifies authorization and sources, not factual truth. Evidence checks compare stored hashes and locators rather than re-reading external bytes; unnoticed external replacement requires immutable storage or additional integrity checks. Commit-entry policy revalidation does not cancel later-revoked in-flight operations (Supplement S2).

Proposition 1 concerns the declared projections, opaque handles, and failure model. Alloy checks finite scopes; selected projections establish bounded correspondence. The formal no-op, schema-uniqueness, and trace-prefix differences are stated in \cref{sec:conformance}. Distributed revocation, other database engines, deletion/redaction, and retention expiry require additional design and evidence.

E1 reuses seven distinct proposal values in constructed cases; the journals are same-study implementations. E2 uses a local scripted fixture, fixed principals, and separate session/fact databases. Its evidence does not extend to human attention, hostile hosts, or cross-database crash atomicity. Historical hosted-agent outcomes and author-involved annotation have separate scoring, missingness, and failure records in Supplement S5.

E3 characterizes one machine with uncontrolled background activity. The admission implementations differ in schema, ORM, validation, features, and connection lifecycle; the materialization timer omits validation/planning. The paired trace study holds the verifier fixed, identifying an access-pattern difference rather than total cost independent of history. Complete workload and validity details are in Supplements S4, S6, and S7.

\subsection{Applying the relation}

Adoption begins with the approval policy and one correction traced through persisted records. Identify the reviewed candidate, authorized value, predecessor check, and changed/unchanged source links. Retain a legal control and a failure probe for each boundary. Supplement S3 provides the worked example and inspection checklist. An independently implemented deployment with operational reviewers would provide the next test of transferability.
